\chapter{The Two Technologies}
\label{ch:methods}

This chapter specifies the two technologies named in Section~\ref{sec:problem}. The first is the predictive controller that moves the three valves and the carbon-dioxide injector. The second is the twin that estimates the model inside that controller. Both are written so that a programmer can implement them, and so that the numerical choices of the case study are the same choices the text describes.

\section{Three-valve nonlinear model predictive control}
\label{sec:nmpc}

\subsection{Time scales and the model that is predicted}

The controller runs every $\Delta t_c = \SI{4}{\second}$. Between samples the command is held. Inside a prediction the hydraulics are algebraic and the composition is integrated with a step of \SI{2}{\second}. The justification is the bulk-modulus argument of \eqref{eq:bulk}. Pressure equalises in a small fraction of a second. Composition equalises on
\begin{equation}
\label{eq:tau}
\tau = \frac{V}{Q_f}
=
\frac{0.080}{20/3600}
=
\SI{14.4}{\second}
\end{equation}
at the design flow. A horizon of \SI{12}{\second} therefore sees most of one residence time and does not pretend to see the next brand. The analyser delay is \SI{12}{\second} as well, so a controller that used only the outlet analysis would be aiming at a beer that has already left. The controller instead integrates the measured inlet flows through \eqref{eq:ode} and \eqref{eq:co2bal}, and it uses the outlet analysis as a slow correction when it arrives. The state that enters the prediction is that integrated composition, plus a small allowance for flow-measurement noise.

The mother-beer analysis enters with its own delay of \SI{8}{\second}. During those eight seconds the predictor still uses the previous mother beer. The error is visible and short. It is not hidden inside a valve gain.

\subsection{Hydraulics inside each shooting node}

At a candidate lift $(l_h,l_w,l_f)$ and a candidate dose, the mix pressure $p_m$ is the unknown. Inlet densities are evaluated at the inlet compositions. The product density is evaluated at the current mixer composition. With the equal-percentage law \eqref{eq:eqpct} and the sizing law \eqref{eq:kv} at $F_P=F_R=1$,
\begin{align}
\label{eq:qhw}
Q_h
&=
K_{v,h}\,R^{\,l_h-1}
\sqrt{\frac{p_h + b_h - p_m}{\rho_h/1000}},
\\
Q_w
&=
K_{v,w}\,R^{\,l_w-1}
\sqrt{\frac{p_w - p_m}{\rho_w/1000}},
\end{align}
when the drops are positive, and zero otherwise. The symbol $b_h$ is the header bias the model believes. On the true plant it is zero. In a nominal, uncommissioned model it is whatever the transmitter adds. The product flow demanded by mass is
\begin{equation}
\label{eq:qfmass}
Q_f^{\mathrm{mass}}
=
\frac{\rho_h Q_h + \rho_w Q_w}{\rho_f}.
\end{equation}
The blend line then drops by \eqref{eq:darcy}, using the Haaland friction factor at the velocity $Q_f^{\mathrm{mass}}/A_f$, and the final valve passes
\begin{equation}
\label{eq:qfvalve}
Q_f^{\mathrm{valve}}
=
K_{v,f}\,R^{\,l_f-1}
\sqrt{\frac{p_m - \Delta p_{\mathrm{line}} - p_b}{\rho_f/1000}}.
\end{equation}
The scalar residual $Q_f^{\mathrm{valve}} - Q_f^{\mathrm{mass}}$ changes sign between a mix pressure just above the tank and a mix pressure just below the lower header. A bracketed scalar solve, Brent's method in the case study, returns $p_m$. Vena-contracta pressures follow from \eqref{eq:pvc} with the recovery factors of the three trims. If the bracket fails, the node is assigned a large cost and the solver steps away from it. That is a numerical guard, not a physical operating point. At the lager solution of Chapter~\ref{ch:estimate} the residual is closed and the cost of the steady problem is numerically zero.

Density is the affine map identified for \SI{2}{\celsius} in the window of this skid,
\begin{equation}
\label{eq:rho}
\rho
=
\big(999.94 - 1.62\,A_w + 4.05\,E_r\big)
\big(1 - \gamma A_w E_r\big).
\end{equation}
The contraction $\gamma$ is zero in the nominal model and is estimated by the twin. Equilibrium pressure uses \eqref{eq:zeta} with the scale $\kappa_{\zeta}$. Alcohol uses \eqref{eq:abv}. Original extract uses \eqref{eq:balling}. None of these inversions is iterative beyond the scalar pressure solve, because $A_w$ and $E_r$ are states, not outputs that have to be searched for.

\subsection{Composition prediction and carbon-dioxide loss}

From a solved hydraulic node the composition is stepped by the forward Euler form of \eqref{eq:ode} and \eqref{eq:co2bal}. The eigenvalue of the mixer is $-1/\tau \approx \SI{-0.07}{\per\second}$. A step of \SI{2}{\second} multiplies a composition error by about $0.87$ per step and does not oscillate. Euler is therefore enough. A production code can substitute a single exponential on the linear mixer without touching the optimiser.

Delivered carbon dioxide is not identical to the mixer concentration. If the final vena contracta lies below equilibrium, a fraction of the supersaturation leaves the liquid in one pass,
\begin{equation}
\label{eq:strip}
c_{\mathrm{out}}
=
c - \eta\,\max\!\big(0,\; c - \kappa_{\zeta}\zeta\,p_{\mathrm{vc},f}\big),
\end{equation}
with $\eta = 0.35$ for this skid. The same test, with a coefficient of \SI{30}{ppb\per\bar}, adds dissolved oxygen when the margin is negative. Both expressions are inactive when the inequalities in \eqref{eq:ineq} hold, which they do at the design point and throughout the closed-loop runs of Chapter~\ref{ch:estimate}. They exist so that a controller is charged for a flashing move in the prediction, rather than discovering the flash in the bright tank.

\subsection{The program solved at each sample}

The horizon is two move blocks of three integration steps. The decision vector is eight numbers,
\begin{equation}
\mathbf{z} = (l_h^{(1)}, l_w^{(1)}, l_f^{(1)}, \dot m^{(1)},\;
l_h^{(2)}, l_w^{(2)}, l_f^{(2)}, \dot m^{(2)}).
\end{equation}
Inside a block the command is constant, and it is still limited by $\dot l_{\max}$ relative to the previous stem position, so a block cannot teleport a valve. The cost is the sum of \eqref{eq:stage} over the six steps. Soft pressure penalties are quadratic outside the margin, with a weight of $6$ on each squared bar, which is large next to a unit tolerance cost. Valve travel is penalised as $0.4\,\|\Delta\mathbf{l}\|^{2}/0.01$, so a simultaneous move of $0.1$ lift on one valve costs about $0.4$. That is less than one tolerance of extract and enough to suppress chatter.

Bounds are the lift limits and a dose limit of \SI{15}{\gram\per\second}. The initial guess is the move that was applied at the previous sample, copied into both blocks. A sequential least-squares quadratic program then reduces the cost. In the case study the iteration limit is six, which is short on purpose: the warm start is already near a good move, and a valve controller that thinks for longer than a fraction of its sample has missed the sample. The production transcription of the same shooting problem is the real-time iteration \citep{diehl2005,gros2020} solved by an interior-point code \citep{wachter2006,andersson2019}. The case-study solver is a stand-in that minimises the same function. It is not a claim that six sequential-quadratic iterations are a certified real-time iteration.

Only the first block is applied. The stem command is limited once more by $\dot l_{\max}\Delta t_c = 0.16$ lift per sample. The dose is applied as a flow setpoint to the injector, which is fast relative to the mixer.

\begin{algorithm}[htbp]
\caption{One sample of three-valve nonlinear model predictive control}
\label{alg:nmpc}
\begin{algorithmic}[1]
\Require Mixer state, delayed mother-beer analysis, header pressures, previous command, brand setpoint
\State Build the initial guess by repeating the previous command in both blocks
\State Predict the horizon from $\mathbf{z}$ through the algebraic network and the mixer step
\State Minimise the summed stage cost subject to lift and dose bounds
\State Rate-limit the first block against the previous stem position
\State If a deadband estimate is active, lead the water-valve command by that band in the direction of travel
\State Apply the three lifts and the dose; hold them until the next sample
\end{algorithmic}
\end{algorithm}

The deadband lead in line 5 is supplied by the twin and is zero until a stick has been observed. The predictive controller of the first technology does not invent it.

\subsection{Steady solution and the relative-gain array}

The same cost, with the derivatives removed and the composition replaced by its algebraic blend, is a four-variable steady program. At the lager brand, on the plant model, the solution is
\begin{equation}
\label{eq:usteady}
l_h = 0.704,\quad
l_w = 0.711,\quad
l_f = 0.840,\quad
\dot m_{\mathrm{CO_2}} = \SI{7.507}{\gram\per\second}.
\end{equation}
The flows are \SI{13.962}, \SI{6.055} and \SI{20.000}{\cubic\metre\per\hour}. The velocities are \SI{1.134}, \SI{0.762} and \SI{1.078}{\metre\per\second} in the \SI{66}{\milli\metre}, \SI{53}{\milli\metre} and \SI{81}{\milli\metre} bores. Mix pressure is \SI{2.855}{bar} absolute against an equilibrium pressure of \SI{1.905}{bar}. The final vena contracta is \SI{2.279}{bar}. Original extract, alcohol and carbon dioxide meet the brand to three decimals, and dissolved oxygen is \SI{9.22}{ppb}, which is the flow-weighted inlet oxygen and no more. The steady cost is below $10^{-8}$. This is the point a grade change leaves and, for the light brand, the kind of point it has to reach.

Linearising the map $(l_h,l_w,l_f)\mapsto(Q_h,Q_w,p_m)$ with a step of $0.012$ lift gives
\begin{equation}
\label{eq:G}
G
=
\begin{bmatrix}
48.05 & -3.21 & 7.64 \\
-2.59 & 23.09 & 2.64 \\
2.09 & 0.89 & -2.14
\end{bmatrix},
\end{equation}
with flows in cubic metres per hour and pressure in bar, per unit lift. The relative-gain array \eqref{eq:rga} is
\begin{equation}
\label{eq:RGAnum}
\Lambda
=
\begin{bmatrix}
0.87 & 0.00 & 0.13 \\
0.00 & 0.95 & 0.05 \\
0.13 & 0.05 & 0.82
\end{bmatrix}.
\end{equation}
The diagonal dominates. Bristol's test supports pairing beer lift with beer flow, water lift with water flow, and final lift with mix pressure. It does not say the off-diagonal terms are zero. Thirteen percent of the beer-flow relative gain sits on the final valve, and the beer valve moves the pressure. A well-tuned decentralised station can hold this point. A simultaneous grade change excites the off-diagonal terms while the mixer is still full of the previous beer. That is the regime in which the predictive controller has work to do, and it is the regime of the first case in Chapter~\ref{ch:estimate}.

The sign of $G_{33}$ is the operator's sign: opening the final valve lowers the mix pressure. The pressure loop of the baseline is therefore reverse-acting. The predictive controller never needs that sentence, because the final lift is just a coordinate of $\mathbf{z}$.

\subsection{The ratio station used as the baseline}

The baseline is the controller the relative-gain array recommends, built so that the comparison is against a competent station rather than against an untuned one.

Given the delayed mother beer, the station inverts \eqref{eq:balling} for the mass fraction $w_{\mathrm{OE}}$ that hits the extract target, and it inverts \eqref{eq:abv} with the linear part of \eqref{eq:rho} for the mass fraction $w_{\mathrm{ABV}}$ that hits the alcohol target. The feedforward fraction is $0.55\,w_{\mathrm{OE}}+0.45\,w_{\mathrm{ABV}}$. A slow integral on the delayed outlet extract, gain $0.006$ per second, trims that fraction. Mass flows are split by $w$ at the demanded product flow, converted to volumetric setpoints with the linear density, and tracked by incremental proportional--integral loops. With $\Delta t_c=\SI{4}{\second}$ the increments are
\begin{align}
\label{eq:pi}
\Delta l_h
&=
0.018\,(e_h-e_h^{-}) + 0.0018\,e_h\,\Delta t_c,
\\
\Delta l_w
&=
0.028\,(e_w-e_w^{-}) + 0.0030\,e_w\,\Delta t_c,
\\
\Delta l_f
&=
0.12\,(e_p-e_p^{-}) + 0.022\,e_p\,\Delta t_c.
\end{align}
Here $e_h$ and $e_w$ are setpoint minus measurement in cubic metres per hour, and $e_p$ is measured mix pressure minus a fixed pressure setpoint. The fixed setpoint is $p_{\mathrm{eq}}+\SI{0.90}{bar}$ evaluated at the lager steady state, which is \SI{2.805}{bar} absolute. The integral gains were chosen so that $K_I\,G_{ii}\,\Delta t_c$ lies near $0.3$ on the diagonal of \eqref{eq:G}. For the beer valve, $0.0018\times 48.05\times 4 \approx 0.35$. For the water valve, $0.0030\times 23.09\times 4 \approx 0.28$. For the pressure loop, $0.022\times 2.14\times 4 \approx 0.19$. These are the SIMC-sized gains for an algebraic process \citep{skogestad2003,astrom2006}. They are not detuned to make the predictive controller look better.

Carbon dioxide is feedforward from \eqref{eq:dosess} using the measured flows and the delayed mother-beer carbonation, plus an incremental correction from the delayed outlet analysis. The same rate limit is applied to the three lifts. Anti-windup is the incremental form itself: the state of the controller is the stem position, and the stem position is clipped to the lift bounds.

What this station cannot do is revise the pressure setpoint when the brand or the mother beer changes $p_{\mathrm{eq}}$, trade extract against alcohol at the tolerance weights when both cannot be met, or refuse a move whose predicted vena contracta flashes. Those are the jobs of \eqref{eq:stage} and \eqref{eq:ineq}.

\subsection{Offset, feasibility and what the horizon does not prove}

A predictive controller that re-optimises from a measured state removes a constant disturbance that the model can see. It does not remove a constant error in its own gain. If $K_v$ is wrong, the predicted flow of a given lift is wrong, the optimiser asks for the wrong lift, and the outlet integral of a ratio station may still be the faster way home. That is why the first technology, deployed on the nominal model, is not expected to dominate the baseline. The second technology exists to make the model that the first technology trusts.

Feasibility of the inequalities is not guaranteed for every mother beer. When extract and alcohol are inconsistent, the hard bounds on the lifts remain feasible and the quality targets become a soft compromise through \eqref{eq:stage}. When a header pressure collapses so far that the product flow cannot be passed without entering the vena contracta, the soft pressure penalty outranks the flow penalty and the controller gives up flow. The case study does not include a collapse of that size. The header step in Chapter~\ref{ch:estimate} is \SI{0.45}{bar}, and every controller keeps a positive margin.

A terminal constraint to the steady set, of the kind \citet{mayne2000} use to prove stability, is the correct production addition to Algorithm~\ref{alg:nmpc}. The case study does not implement it. The horizon is short, the rate limit is hard, and the steady program is feasible at both brands. The numerical evidence in Chapter~\ref{ch:estimate} is an estimation of closed-loop error, not a Lyapunov certificate. Saying so keeps the first technology in the class of engineering predictive control reviewed by \citet{qin2003} and \citet{maciejowski2002}, where certificates are added when the model and the constraints have settled.

\subsection{A grade change, computed before the valves move}

The light brand asks for $w=0.58$ rather than $0.70$. Product mass flow at the lager density is about \SI{5.61}{\kilogram\per\second}. The beer mass flow therefore falls from about \SI{3.92}{\kilogram\per\second} to about \SI{3.25}{\kilogram\per\second}, and the water mass flow rises by the same \SI{0.67}{\kilogram\per\second}. In volumetric units the beer setpoint moves from \SI{13.96}{\cubic\metre\per\hour} toward about \SI{11.6}{\cubic\metre\per\hour}, and the water setpoint rises from \SI{6.06}{\cubic\metre\per\hour} toward about \SI{8.4}{\cubic\metre\per\hour}. Divided by the diagonal gains, those are lift changes of a few hundredths on the beer valve and about a tenth on the water valve. At \SI{0.04}{lift\per\second} that water-valve move takes about three seconds. Holding the inlet exactly at the new brand after that move would still leave an extract error near \SI{0.7}{\degree}P one residence time later, because $\mathrm{e}^{-1}$ times the \SI{1.93}{\degree}P step has not yet entered the \SI{0.05}{\degree}P band. A controller that is inside the band in \SI{16}{\second} has therefore used an inlet lighter than the new brand while the vessel empties. Section~\ref{sec:firstsample} records that move on the commissioned model. A controller that is still outside the band two minutes later is limited by its model or by its tuning, not by the pipe. Chapter~\ref{ch:estimate} uses that comparison directly.

The carbon-dioxide setpoint falls from \SI{5.20}{\gram\per\litre} to \SI{5.00}{\gram\per\litre} at the same moment. Because the water fraction rises, the dose does not fall in proportion: more uncarbonated water has to be made up. The predictive controller sees both effects in \eqref{eq:co2bal}. The baseline sees them in \eqref{eq:dosess}, which is algebraically the same at steady state and is blind to the trajectory between.

\subsection{Installed gain at the lager lifts}
\label{sec:installed}

The diagonal of \eqref{eq:G} is not an arbitrary linearisation. It is the derivative of \eqref{eq:kv} at the lifts of Table~\ref{tab:steady}, with the mix pressure free to move. Holding the drop fixed isolates the valve. At beer lift $0.704$, $\Phi=0.314$ and $\ln R=\ln 50=3.912$, so
\begin{equation}
\label{eq:dphi}
\frac{\mathrm{d}\Phi}{\mathrm{d}l}
=
\ln R\cdot\Phi
=
3.912\times 0.314
=
1.228.
\end{equation}
A fixed drop would therefore scale the beer flow by $\ln R$ for a unit of lift,
\begin{equation}
\left.\frac{\partial Q_h}{\partial l_h}\right|_{\Delta p}
=
Q_h\ln R
=
13.96\times 3.912
=
\SI{54.6}{\cubic\metre\per\hour}.
\end{equation}
The network gain in \eqref{eq:G} is $48.05$, about twelve percent smaller. The missing part is the rise in mix pressure that follows an opening of the beer valve. That rise shrinks the drop the valve was using, so the installed gain is below the fixed-drop gain. The same comparison on the water valve, $Q_w\ln R = 6.055\times 3.912 = 23.7$ against the network value $23.09$, is almost flat, because the water valve is the smaller flow and moves $p_m$ less. The final valve has the opposite sign: opening it lowers $p_m$, and the inlet valves then pass more flow, which is the positive $G_{13}=7.64$ and $G_{23}=2.64$.

This is the installed characteristic of \citet{iec60534,fisher2005} written as a number the controller can use. An equal-percentage trim was chosen so that the local gain stays near $\ln R$ times the current flow over the whole stroke. At the lager point that product is about \SI{55}{\cubic\metre\per\hour} per unit lift on the beer valve before the network correction, and the correction is modest. A linear trim of the same $K_v$ would have put most of its gain near the shut end and very little at lift $0.70$, which is exactly where this skid lives. The predictive controller does not assume the gain is constant. It re-solves \eqref{eq:kv} at every node. The linear gain is used to size the baseline in \eqref{eq:pi} and to interpret the relative-gain array, not to predict the brand.

\subsection{Bristol's array from the measured gain}
\label{sec:bristol}

The relative-gain array is $G$ multiplied elementwise by the transpose of $G^{-1}$ \citep{bristol1966}. The point of writing the arithmetic is to show that the pairing in Section~\ref{sec:problem} is a property of these three numbers and not a convention copied from a flow--flow--pressure textbook.

Denote the rows of \eqref{eq:G} by the beer flow, the water flow and the mix pressure. The determinant is dominated by the product of the diagonal, $48.05\times 23.09\times(-2.14)$, which is about $-2370$, and the off-diagonal corrections do not change its sign. The $(1,1)$ entry of $G^{-1}$ is the beer-flow response to a beer-lift move after the other two outputs have been held at zero by the other two inputs. Multiplied back by $G_{11}$ it is $\Lambda_{11}=0.87$. The same construction gives $\Lambda_{22}=0.95$ and $\Lambda_{33}=0.82$. Each diagonal entry is the fraction of the open-loop gain that survives when the other loops are closed perfectly. A value near one means the other loops do not steal the actuator. A value near zero would mean the recommended pairing is the wrong one. A negative value would mean the loops fight.

Thirteen percent of the beer-flow relative gain sits on the final valve, and five percent of the water-flow relative gain sits there too. Those are the terms a grade change excites, because a grade change moves all three lifts in the same few seconds, and perfect closure of the other loops is exactly what has not happened yet. At steady state, after the ratio station has settled, the same thirteen percent is a small trim on the pressure loop. The array therefore explains both observations in Chapter~\ref{ch:estimate}: the ratio station can hold a brand, and it is slow when the brand changes. It does not by itself choose the nonlinear map. Balling's denominator is outside $G$, and no arrangement of these nine gains puts it back.

\subsection{How long the horizon has to be}
\label{sec:horizon}

The mixer time constant at the design flow is $\tau=\SI{14.4}{\second}$. The control sample is $\Delta t_c=\SI{4}{\second}$. The prediction uses two move blocks and three Euler steps of \SI{2}{\second} inside each block, so the horizon is
\begin{equation}
T_H = 2\times 3\times 2 = \SI{12}{\second},
\end{equation}
which is $0.83\,\tau$. A step in the inlet composition has, by \eqref{eq:exactdisc}, reached a fraction $1-\mathrm{e}^{-T_H/\tau}=1-\mathrm{e}^{-0.83}=0.56$ of its final value at the end of the horizon. The controller therefore sees a little more than half of the quality move that a valve decision will eventually produce. It does not see the whole settling tail. The tail is left to the next sample, which is the receding-horizon argument: only the first block is applied, and the remainder of the response is re-predicted when the state has moved.

Lengthening the horizon to $2\tau$, about \SI{29}{\second}, would show $86$ percent of the step and would multiply the number of network solves by about two and a half. On a skid whose disturbance is a brand change that is known at the sample it begins, the extra tail is a smaller extract error at the far nodes, already shrinking by $\mathrm{e}^{-\alpha}$ each step, and the decision that matters is the first block. The case study keeps $T_H=\SI{12}{\second}$ for that reason. A disturbance that arrived as an unmeasured ramp over a minute, for example a slow header decay, would justify the longer horizon. The header step in the second scenario is a step, and \SI{12}{\second} covers its onset.

The two blocks are not two independent valve trajectories. Inside a block the command is held, so the eight-vector $\mathbf{z}$ is
\begin{equation}
\mathbf{z}
=
\big(l_h^{(1)},l_w^{(1)},l_f^{(1)},\dot m_{\mathrm{CO_2}}^{(1)},
l_h^{(2)},l_w^{(2)},l_f^{(2)},\dot m_{\mathrm{CO_2}}^{(2)}\big).
\end{equation}
Six Euler steps then see only two values of each actuator. The rate limit is enforced between the previous applied command and the first block, and again between the first block and the second. With a limit of $0.04$ lift per second, the first block may differ from the applied command by at most $0.04\times 4=0.16$ lift, and the second block may differ from the first by the same amount. In two samples the water valve can travel $0.32$ lift. The grade-change arithmetic of the previous subsection needs only about a tenth of a lift on that valve. The rate limit is not what makes the vessel take \SI{14}{\second}. The volume is.

\subsection{A move that the rate limit does price}
\label{sec:travel}

Travel in Table~\ref{tab:kpi} is the sum of absolute lift increments over the run. On the grade change the ratio station travels $7.09$ and both predictive controllers travel about $1.65$. The difference is not a smoother optimum in the abstract. It is the incremental law \eqref{eq:pi} reacting to a flow error that the mixer has not yet produced. At the sample of the setpoint change the ratio station computes new flow setpoints immediately, the flows are still the lager flows, and the errors $e_h$ and $e_w$ are the full step of about $2.4$ and $2.3$ cubic metres per hour. The proportional increments are $0.018\times 2.4$ and $0.028\times 2.3$, a few hundredths of lift, and the integral keeps adding while the vessel is still exchanging. Each later sample sees a flow error of the opposite sign as the valves overshoot the new ratio, and the stems reverse. The sum of those reversals is the travel of $7$.

The predictive controller applies one block that already points at the new steady lifts, then a second block that is close to the first, and then stops. The stage cost charges $0.4\,\|\Delta l\|^{2}/0.01$ for that one move and charges nothing further once the prediction says the vessel will arrive. The travel of $1.65$ is that move plus the small trims that follow the analyser. It is larger than the one-tenth lift of the steady-state arithmetic because the final valve and the dose move as well, and because the first sample cannot jump the whole way when the warm start is the lager command and the iteration budget is six. The budget spends itself on the first sample of the grade change. Later samples start next to a solution and barely move.

The same accounting explains the third scenario, where the commissioned controller travels $0.63$ and the ratio station travels $0.20$. The extra travel is the deadband lead. The stem is commanded across a band of width $0.03$ that produces no flow, so that the flow, when it starts, is the flow the brand asked for. A travel column that is read as ``smaller is always better'' would punish the only controller that has understood the valve. The column is reported so that a reversal hunt, of the kind the untuned absolute-form law produced before the gains of \eqref{eq:pi} were fixed, cannot hide inside a good extract integral.

\subsection{Where the carbon-dioxide dose sits in the same vector}
\label{sec:dosegain}

The dose is the fourth coordinate of each block and it does not appear in $G$, because $G$ was identified on the three lifts with the dose doing whatever holds \SI{5.20}{\gram\per\litre}. Its own gain is algebraic. One gram per second of carbon dioxide, divided by \SI{20}{\cubic\metre\per\hour}, is
\begin{equation}
\frac{1\,\mathrm{g\,s^{-1}}\times 3600}{20\times 10^{3}\,\mathrm{L\,h^{-1}}}
=
\SI{0.180}{\gram\per\litre}
\end{equation}
at the design flow, before the mixer lags it. The stage cost treats \SI{0.05}{\gram\per\litre} as one tolerance, so a dose error of $0.05/0.180=\SI{0.28}{\gram\per\second}$ costs the same, at steady state, as one tolerance of original extract. The optimiser can therefore spend a fraction of a gram per second to protect carbonation without giving up the brand, and it cannot hide a one-gram-per-second mistake. That scaling is why the dose belongs in $\mathbf{z}$ rather than in a slave loop with an unrelated gain. A slave loop can still implement the command. The predictive program decides it, jointly with the lifts, because a water-valve move changes the uncarbonated inlet at the same moment the dose should change.

At the lager point the required dose is \SI{7.507}{\gram\per\second}. The upper bound in \eqref{eq:ineq} is \SI{15}{\gram\per\second}, which is twice the design and is the injector size used in the study. The bound is inactive at both brands. It would become active if the mother beer arrived nearly flat and the product specification stayed at \SI{5.2}{\gram\per\litre}, a case the commissioning range does not include. The program already has the inequality. No new term is required.

The equilibrium pressure moves when the dose moves, because $c$ moves and $p_{\mathrm{eq}}=c/(\kappa_\zeta\zeta)$. A dose increase of \SI{1}{\gram\per\second} raises $c$ by about \SI{0.18}{\gram\per\litre} and $p_{\mathrm{eq}}$ by about $0.18/2.73=\SI{0.066}{bar}$ once the mixer has caught up. The margin constraint sees that rise. A controller that carbonates in a slave loop and only afterwards checks the valve throat will discover the lost margin one analyser delay later. The joint program discovers it inside the shooting node, while the dose is still a decision.

\subsection{The first sample of a grade change}
\label{sec:firstsample}

Table~\ref{tab:horizon} collects the numbers that fix the first sample, so that the horizon is a design choice with units rather than a list of solver options.

\begin{table}[htbp]
\centering
\caption{Prediction budget at the lager flow of \SI{20}{\cubic\metre\per\hour}.}
\label{tab:horizon}
\begin{tabular}{lr}\toprule
Quantity & Value \\\midrule
Mixer volume & \SI{0.080}{\cubic\metre} \\
Residence time $\tau$ & \SI{14.4}{\second} \\
Control sample $\Delta t_c$ & \SI{4}{\second} \\
Euler step inside a block & \SI{2}{\second} \\
Move blocks $\times$ steps & $2\times 3$ \\
Horizon $T_H$ & \SI{12}{\second} \\
Fraction of a step seen, $1-\mathrm{e}^{-T_H/\tau}$ & 0.56 \\
Decision length & 8 \\
Iteration limit & 6 \\
Lift rate limit & \SI{0.04}{lift\per\second} \\
Largest legal first move & 0.16 lift \\
\bottomrule\end{tabular}
\end{table}

At the sample the light brand is requested, the vessel still holds the lager. Original extract is \SI{11.61}{\degree}P against a new target of \SI{9.687}{\degree}P. The error is \SI{1.93}{\degree}P, which is $38.5$ tolerances, and one stage of \eqref{eq:stage} charges about $1.5\times 10^{3}$ on extract alone. Six stages, if the command were left at the lager lifts, charge about $10^{4}$. The horizon sees only $56$ percent of a residence-time step, so a move that merely places the inlet on the new brand produces a predicted extract error that is still enormous at the last node. The solver spends the rate budget to make the inlet lighter than the brand, because that is the only way the predicted error inside \SI{12}{\second} falls by a useful fraction.

On the commissioned model the plant is stepped and then the controller runs, so the new command appears one second later. The lager lifts are $0.703$, $0.713$ and $0.835$. The next applied lifts are $0.543$, $0.873$ and $0.675$. Each coordinate has moved by $0.160$, which is the rate limit of $0.04\times 4$ and not the tenth of a lift that separates the two steady states. Beer flow falls from \SI{13.8}{\cubic\metre\per\hour} to \SI{6.4}{\cubic\metre\per\hour} and water flow rises from \SI{6.1}{\cubic\metre\per\hour} to \SI{10.2}{\cubic\metre\per\hour}. The inlet mass fraction is near $0.39$, not $0.58$. The dose, which had been \SI{7.5}{\gram\per\second}, is raised with the extra water and reaches about \SI{11.6}{\gram\per\second} on the following sample. Four seconds after the upset the extract error is \SI{0.87}{\degree}P. A pure decay at $\tau=\SI{14.4}{\second}$ would still be at $\mathrm{e}^{-4/14.4}\times 1.93=\SI{1.46}{\degree}P$. The difference is the light inlet.

The next samples spend the budget in the other direction. By \SI{16}{\second} the lifts are near $0.66$, $0.81$ and $0.79$, the flows are near \SI{11.4}{\cubic\metre\per\hour} and \SI{8.7}{\cubic\metre\per\hour}, and the extract error is \SI{0.037}{\degree}P, which is the first sample of the strict band in Table~\ref{tab:kpi}. One intermediate sample cuts the product flow to \SI{13.2}{\cubic\metre\per\hour} while the ratio is being restored. That sample is visible in the flow integral. Of the \SI{51}{(\cubic\metre\per\hour)\cdot s} charged to this controller, almost all of it sits in those first \SI{16}{\second}. The extract integral over the whole run is \SI{10.75}{\degree}P$\cdot$s. The same step, decayed at the residence time from an inlet that jumped exactly onto the new brand and integrated to infinity, would cost
\begin{equation}
\int_{0}^{\infty} 1.93\,\mathrm{e}^{-t/14.4}\,\mathrm{d}t
=
1.93\times 14.4
=
27.8
\end{equation}
in degree-Plato seconds. The commissioned run is below that number because the inlet was lighter than the brand, and the flow weight in \eqref{eq:stage} is what the optimiser paid in exchange. The payment is finite. The two other stations on the same run finish at \SI{31.30}{\degree}P$\cdot$s and \SI{32.29}{\degree}P$\cdot$s, and neither enters the band. Their integrals contain the vessel and a model, or a tuning, that keeps feeding the wrong inlet after the vessel could have been emptied.

The nominal predictive controller is the second of those two numbers. It uses the same horizon, the same weights and the same rate limit. The only difference on that run is $\theta$. An extract integral that matches the ratio station, and a tail of about \SI{0.19}{\degree}P, is what an uncommissioned capacity and a \SI{0.10}{bar} header bias produce. The horizon arithmetic does not repair it. The twin does. The six-iteration budget is also visible: the intermediate flow dip is a correction a converged program would have spread more evenly. It does not keep the extract outside the band, and it is one reason the production solver of Section~\ref{sec:rti} replaces the iteration limit by a single warm-started quadratic program.


\section{Physics-informed pressure digital twin}
\label{sec:twin}

\subsection{What the twin contains}

The twin is a second copy of the equations of Section~\ref{sec:nmpc}, with a parameter vector that the skid is allowed to disagree about, a residual correction that is allowed to be zero, and an interface that hands the commissioned copy to Algorithm~\ref{alg:nmpc}. In the classification of \citet{kritzinger2018} it is a twin because that hand-off moves the valves. A display of the same equations would be a shadow.

The parameter vector is
\begin{equation}
\label{eq:theta}
\theta
=
\big(
\log K_{v,h},\;
\log K_{v,w},\;
\log K_{v,f},\;
\log \varepsilon,\;
\log K,\;
b_h,\;
\kappa_{\zeta},\;
\gamma
\big).
\end{equation}
Logarithms keep the capacities, the roughness and the minor-loss coefficient positive. The bias $b_h$ is in bar and may have either sign. The solubility scale and the contraction are estimated from different experiments than the valve steps, because they are barely visible in those steps. They are stored in $\theta$ because the predictor reads them there.

The nominal twin, before commissioning, is the model a project would start from: capacities \SI{29.4}, \SI{13.2} and \SI{58.0}{\cubic\metre\per\hour} against true values \SI{32}, \SI{12} and \SI{64}; roughness \SI{15}{\micro\metre} against \SI{2}{\micro\metre}; minor loss $8$ against $14$; solubility scale $1$ against $0.93$; contraction $0$; and a belief that the beer-header transmitter is truth when it reads \SI{0.10}{bar} high. Those errors are the size of a worn or mis-sized trim and a transmitter that has not been zeroed. They are not a caricature.

\subsection{Commissioning experiment}

Twenty-eight steady points are collected. Each point holds a random lift triple in $[0.40,0.82]$ and a random dose, long enough for the mixer to settle, and records the three flows and the mix pressure. Flow noise is \SI{0.03}{\cubic\metre\per\hour} standard deviation. Pressure noise is \SI{0.008}{bar}. The beer-header measurement used by the estimator is the true header plus \SI{0.10}{bar} plus that noise. Inlet composition is the mother beer of the design point, known because the campaign is run on one tank. This is a few hours of a skid, not a month of history. The point of a physics-based estimator is that a few hours are enough \citep{zheng2023,kadlec2009}.

The hydraulic residual stacks four normalised errors per point,
\begin{equation}
\label{eq:resid}
r
=
\begin{bmatrix}
(Q_h-\tilde Q_h)/0.08 \\
(Q_w-\tilde Q_w)/0.05 \\
(Q_f-\tilde Q_f)/0.08 \\
(p_m-\tilde p_m)/0.02
\end{bmatrix},
\end{equation}
and $\theta$ enters $Q$ and $p_m$ through the network of Section~\ref{sec:nmpc}. The estimator minimises $\|r(\theta)\|^{2}$ by a bound-aware nonlinear least-squares method, starting from the nominal capacities. An approximate covariance follows from the Jacobian $J$ of the converged residual,
\begin{equation}
\label{eq:cov}
\mathrm{Cov}(\hat\theta)
\approx
\hat\sigma^{2}\,(J^{\mathsf T}J)^{-1},
\qquad
\hat\sigma^{2}
=
\frac{\|r\|^{2}}{N-n_{\theta}}.
\end{equation}
Standard errors on $K_v$ are the standard errors on $\log K_v$, multiplied by $K_v$. This is the local Gaussian approximation used throughout nonlinear least squares. It is honest about precision only when the parameters are locally identifiable. Chapter~\ref{ch:estimate} shows that for roughness and minor loss they are not, and the enormous standard errors are the estimator saying so.

\subsection{Parameters that need their own experiment}

Solubility is estimated from one bench point. A sample of known $A_w$, $E_r$ and dissolved carbon dioxide is brought to equilibrium at \SI{2}{\celsius}, the headspace pressure is read, and
\begin{equation}
\label{eq:kappahat}
\hat\kappa_{\zeta}
=
\frac{c}{\zeta_{\mathrm{poly}}(T,E,\mathrm{ABV})\,p_{\mathrm{meas}}}.
\end{equation}
One point identifies one scale. It does not re-identify the polynomial. A brewery that changes beer style enough to leave the polynomial's range repeats the bench point on the new style. That is a laboratory sample, not a flashing test on the line.

Contraction is estimated from the Coriolis density. At each steady point,
\begin{equation}
\label{eq:gammahat}
\hat\gamma
=
\mathrm{median}_{i}
\left(
\frac{1 - \tilde\rho_i/\rho_{\mathrm{lin}}(A_{w,i},E_{r,i})}{A_{w,i}E_{r,i}}
\right),
\end{equation}
taken over points where $A_w E_r > 1$ so that the denominator is not a trace of noise. The median rejects a single bad density reading. The linear map $\rho_{\mathrm{lin}}$ is the parenthesis in \eqref{eq:rho}.

Effective line resistance is not a seventh search variable. After $\varepsilon$ and $K$ have been estimated, or have failed to be estimated, it is computed from \eqref{eq:closs} at $\mathrm{Re}=4\times 10^{4}$, which is the blend-line Reynolds number at the design point. The twin may be rebuilt with that single $C_L$ in place of the pair $(\varepsilon,K)$. In the case study the deployed twin keeps the estimated pair because, although each member is nonsense, the sum $C_L$ lands within three percent of the truth. The text reports both facts so that the next skid estimates $C_L$ directly and does not publish a roughness.

\subsection{The residual network and the mass-balance gate}

After the physical parameters are fixed, the flow residual on the training points is a target for a small network. The input is the standardised vector $(l_h,l_w,l_f,A_w,E_r,p_m)$. The network has one hidden layer of eight hyperbolic-tangent units and a linear output of three flow corrections. Twenty-two of the twenty-eight points train it. Six are held out. Training is Adam for 900 steps at learning rate $0.02$, with a weight decay of $2\times 10^{-3}$.

Every fifth step adds a physics penalty. If $q$ is the physically predicted flow plus the network correction, the penalty is the squared mass imbalance
\begin{equation}
\label{eq:masspen}
\big(\rho_h q_h + \rho_w q_w - \rho_f q_f\big)^{2},
\end{equation}
scaled so that it does not drown a residual that is already at the sensor floor. This is the blend-skid form of the collocation penalty of \citet{raissi2019}. The conservation law is algebraic rather than a differential equation, because the quantity being corrected is a steady flow. \citet{cai2021} make the same substitution whenever a fluid constraint is a flux balance rather than a momentum residual.

The gate is the held-out root-mean-square flow error. If the network increases it, the deployed twin uses $\Delta Q = 0$. The architecture stays in the software so that a later campaign, on a valve whose characteristic is not equal-percentage, can open the gate without a rewrite. \citet{wu2023cerd} and \citet{zheng2023} deploy the network because their reactors have a kinetic discrepancy that the parameters do not absorb. This skid, after $K_v$ and $b_h$ are estimated, does not. The gate is how the twin tells those two situations apart.

\begin{algorithm}[htbp]
\caption{Commissioning of the physics-informed pressure twin}
\label{alg:twin}
\begin{algorithmic}[1]
\Require Step-test records, one bench equilibrium pressure, Coriolis densities
\State Estimate $\log K_v$, $\log\varepsilon$, $\log K$ and $b_h$ by least squares on \eqref{eq:resid}
\State Compute standard errors from \eqref{eq:cov}; if a standard error exceeds the estimate, replace that parameter by the identifiable combination $C_L$
\State Estimate $\kappa_{\zeta}$ from \eqref{eq:kappahat} and $\gamma$ from \eqref{eq:gammahat}
\State Train the residual network on all but six step tests, with the penalty \eqref{eq:masspen}
\State If the held-out flow error rises, set the residual to zero
\State Publish $\theta$ and the residual gate to the predictive controller
\end{algorithmic}
\end{algorithm}

\subsection{The fast state, and the deadband}

During a run the slow parameters stay at the commissioned values. The mixer state is integrated from the measured flows and the delayed mother-beer analysis. This observer is a soft sensor in the sense of \citet{kadlec2009}, and it is almost a balance rather than a regression: if the flows and the inlet composition are right, the mixer composition is determined. Noise of a hundredth of a percent by weight is added in the case study to represent the imperfection of that statement. An unscented filter \citep{julier2004} or a moving-horizon estimator \citep{rao2003} is the right replacement when an inlet meter drops out and the outlet analysis has to carry the state alone. With all three Coriolis meters healthy, the balance is the estimator.

Water-valve stiction is detected as a deadband $S$: the stem does not move until the command has led the reported position by $S$. The case-study fault is $S=0.03$ lift, with no slip jump, which is the stick band of \citet{choudhury2005} reduced to one parameter. The inverse used by the twin-backed controller is
\begin{equation}
\label{eq:invstick}
l_{w,\mathrm{cmd}}
=
l_{w,\mathrm{des}} + S\,\mathrm{sign}(l_{w,\mathrm{des}} - l_{w,\mathrm{pos}}),
\end{equation}
applied only when the desired move exceeds a small threshold, so that a valve already at the setpoint is not dithered. \citet{jelali2010} discuss the limit cycle that a naive inverse can excite. The threshold and the rate limit are what keep \eqref{eq:invstick} from becoming a knocker. The predictive controller that does not receive $S$ has no inverse. It sees a flow that refuses small moves and, if its gain is wrong as well, it walks away from the setpoint.

\subsection{How the twin changes the first technology}

Algorithm~\ref{alg:nmpc} does not contain $\theta$. It contains a model. Commissioning replaces that model. Three substitutions are visible in the closed loop of Chapter~\ref{ch:estimate}.

The capacities and the header bias are the substitutions that repair the grade change. An eight-percent error in $K_{v,h}$ is an error of about \SI{1}{\cubic\metre\per\hour} at the lager point, because the valve is not fully open and the equal-percentage gain is steep. The predictive controller converts that flow error into a composition error and then holds it, because the model believes the composition is already right. The ratio station, which integrates measured flow rather than predicted flow, is less sensitive to $K_v$ and more sensitive to the interaction during the move. The numerical comparison falls out of that difference.

The solubility scale and the contraction are the substitutions that keep the pressure constraint and the alcohol conversion honest. At the design point $\kappa_{\zeta}=0.93$ lowers $\zeta$ from the bare polynomial by seven percent and raises $p_{\mathrm{eq}}$ by the same fraction. A controller that ignores the scale guards the wrong pressure. The contraction shifts density by about \SI{4}{\kilogram\per\cubic\metre} on the mother beer and by about \SI{2}{\kilogram\per\cubic\metre} on the lager. Alcohol by volume, through \eqref{eq:abv}, moves in proportion. The shift is smaller than the alcohol tolerance and is still worth removing, because the tolerance budget should be spent on the mother beer, not on the equation of state.

The deadband inverse is the substitution that matters when the water valve sticks. It does not change the steady mathematics of $w$. It changes the first samples after a small flow error appears.

The residual network, in this campaign, is the substitution that is not made. Publishing that fact is a duty of the estimator. A twin that always adds its network, after the network has failed a hold-out test, is a softer sensor with a harder name.

\subsection{Information, interfaces and the safety boundary}

The twin reads flows, pressures, stem positions, temperature and quality through the OPC UA model of the skid \citep{iec62541}. It writes a recommended $(l_h,l_w,l_f,\dot m_{\mathrm{CO_2}})$ at \SI{4}{\second} intervals. The programmable controller remains the only device that energises a solenoid. It enforces the rate limit again, the lift limits again, and a set of interlocks the optimiser is not allowed to relax: a low bright-tank pressure, a high mixer pressure, a loss of the water-deaeration permissive, and a quality hold if dissolved oxygen or a water analysis crosses the alarm. ISA-88 names the brand change as a phase whose running state may be held by those permissives \citep{isa88}. The predictive controller is a step inside the running state. It is not a replacement for the phase logic.

This boundary is the practical content of the warning in \citet{fuller2020} and \citet{goffi2025} that a twin which can move equipment has to be commissioned inside the interlocks, not beside them. The estimation problem is unaffected. The interlocks define a subset of $\mathbf{u}$ that the least-squares model never sees, because the campaign is run with the interlocks healthy. Extrapolating the twin into a tripped skid is not a use this thesis supports.

\subsection{A hand solution of the lager network}
\label{sec:hand}

The optimiser's steady point can be checked with a calculator, and the check is worth writing down because it is the same arithmetic every shooting node repeats. Take the lifts of \eqref{eq:usteady} and the mother-beer density \SI{1011.9}{\kilogram\per\cubic\metre}. The equal-percentage factor at beer lift $0.704$ is
\begin{equation}
\Phi(0.704) = 50^{-0.296} = \mathrm{e}^{-0.296\ln 50} = \mathrm{e}^{-1.158} = 0.314.
\end{equation}
The effective capacity is $32\times 0.314 = 10.05$ cubic metres per hour. The mix pressure solved by the scalar residual is \SI{2.855}{bar}, so the beer-valve drop is $4.80-2.855=\SI{1.945}{bar}$. With specific gravity $1.012$,
\begin{equation}
Q_h = 10.05\sqrt{1.945/1.012} = 10.05\times 1.386 = \SI{13.93}{\cubic\metre\per\hour}.
\end{equation}
The network solution is \SI{13.962}{\cubic\metre\per\hour}. The difference is the difference between this one-pass density and the density the scalar solver uses after the contraction and the blend have been updated. It is two tenths of a percent. The same arithmetic on the water valve, $\Phi(0.711)=0.323$, effective capacity $3.87$, drop $5.30-2.855=\SI{2.445}{bar}$, specific gravity essentially $1$, gives
\begin{equation}
Q_w = 3.87\sqrt{2.445} = \SI{6.05}{\cubic\metre\per\hour},
\end{equation}
against the solved \SI{6.055}{\cubic\metre\per\hour}.

The blend-line velocity at \SI{20}{\cubic\metre\per\hour} in an \SI{81}{\milli\metre} bore is
\begin{equation}
v_f = \frac{20/3600}{\pi(0.0405)^{2}} = \SI{1.078}{\metre\per\second},
\end{equation}
which is the value in Table~\ref{tab:steady}. The Reynolds number with $\rho\approx\SI{1009}{\kilogram\per\cubic\metre}$ and $\mu=\SI{2.2}{\milli\pascal\second}$ is $4.0\times 10^{4}$. Haaland's formula at $\varepsilon=\SI{2}{\micro\metre}$ returns $f=0.0218$. Then
\begin{equation}
f\frac{L}{D} = 0.0218\times\frac{18}{0.081} = 4.85,
\qquad
C_L = 4.85 + 14 = 18.85,
\end{equation}
and the line drop is
\begin{equation}
\Delta p_{\mathrm{line}} = 18.85\times\frac{1009\times 1.078^{2}}{2}\times 10^{-5} = \SI{0.111}{bar}.
\end{equation}
The pressure immediately upstream of the final valve is therefore about \SI{2.74}{bar}. The final-valve drop down to the tank at \SI{2.40}{bar} is about \SI{0.34}{bar}. With $F_L=0.86$,
\begin{equation}
p_{\mathrm{vc},f} = 2.74 - \frac{0.34}{0.86^{2}} = 2.74 - 0.46 = \SI{2.28}{bar},
\end{equation}
which matches the solved vena contracta of \SI{2.279}{bar}. Equilibrium at \SI{5.20}{\gram\per\litre} and $\zeta=\SI{2.730}{\gram\per\litre\per\bar}$ is \SI{1.905}{bar}. The margin at the narrowest point in the final valve is \SI{0.37}{bar}, and the margin at the mix node is \SI{0.95}{bar}. A designer who stops at the mix-node margin and never computes \eqref{eq:pvc} will believe the skid has a bar of room when the valve throat has a third of that.

The carbon-dioxide check was already done in Section~\ref{sec:build} and closes to \SI{7.51}{\gram\per\second}. The hand solution and the program are the same model. The program is used because a grade change evaluates the residual hundreds of times, not because the residual is obscure.

\subsection{The mixer as a discrete linear system}
\label{sec:discrete}

Once the flows are known, each composition channel is a linear filter. Let $\alpha = Q_f\Delta t/V$ with $\Delta t=\SI{2}{\second}$ and $V=\SI{0.08}{\cubic\metre}$ at the design flow, so $\alpha = 2/14.4 = 0.139$. The exact discretisation of a single well-mixed volume is
\begin{equation}
\label{eq:exactdisc}
x_{k+1} = \mathrm{e}^{-\alpha}\,x_k + \big(1-\mathrm{e}^{-\alpha}\big)\,x_{\mathrm{in}},
\end{equation}
and the Euler step used in the shooting code is the first-order form $x_{k+1}=(1-\alpha)x_k+\alpha x_{\mathrm{in}}$. At this $\alpha$ the two gains on the old state are $\mathrm{e}^{-\alpha}=0.870$ and $1-\alpha=0.861$. The relative difference is one percent. Over a six-step horizon the compounded difference in a step response is a few hundredths of the step, which is smaller than the analyser noise on original extract. Euler is kept because its derivative with respect to the flow, which changes $\alpha$, is obvious to the shooting code. A production build can switch to \eqref{eq:exactdisc} without touching the cost.

The four channels $A_w$, $E_r$, $c$ and $c_{\mathrm{DO}}$ share $\alpha$ and do not share $x_{\mathrm{in}}$. Alcohol and extract see a mass-weighted inlet. Carbon dioxide and oxygen see a volume-weighted inlet plus, for carbon dioxide, the dose divided by $Q_f$. The state transition is therefore diagonal. The output map is not. Original extract and alcohol are nonlinear functions of $A_w$, $E_r$ and $\rho$, and $p_{\mathrm{eq}}$ is a nonlinear function of those and of $c$. That is why a linear model predictive controller built on $G$ in \eqref{eq:G} can coordinate the valves and still miss the brand: $G$ does not contain Balling's denominator. The nonlinear shooting exists for the output map, not because the mixer is a difficult dynamic system.

A linear controller is still the right tool for the diagnosis in \eqref{eq:RGAnum}. It is the wrong tool for the setpoint. The numerical size of the mistake was computed in Section~\ref{sec:problem}: a dilution proportional to original extract misses the true original extract by \SI{0.19}{\degree}P. Any linear gain from flow ratio to original extract that was identified at one brand will carry an offset of that order at the next brand. Integral action removes the offset after the analyser delay. The nonlinear map removes it in the prediction.

\subsection{What one tolerance costs, in the units of the stage cost}
\label{sec:weights}

The denominators in \eqref{eq:stage} are the tolerances, so the numerical weight of each channel is a statement about which miss the brewery will trade. It is worth evaluating the cost at a few concrete errors so the weights are not an unexamined vector.

One tolerance of original extract, with every other error zero, costs $1$. One tolerance of alcohol costs $1$. One tolerance of carbon dioxide costs $1$. A flow error of \SI{0.5}{\cubic\metre\per\hour}, which is \SI{2.5}{\percent} of the design flow, costs $0.35$. A valve move of $0.1$ lift from the previous command costs $0.4\,(0.1)^{2}/0.01 = 0.4$. A mix pressure \SI{0.1}{bar} inside the forbidden side of the margin costs $6\times 0.01 = 0.06$, and \SI{0.5}{bar} inside it costs $1.5$. The pressure penalty therefore overtakes a quality tolerance only when the violation is several tenths of a bar. Small numerical noise in $p_m$ does not dominate the brand. A flashing valve does.

Apply the same arithmetic to the infeasible mother beer of Section~\ref{sec:s2}. The extract-only point misses alcohol by \SI{0.207}{\percent} by volume, which is $0.207/0.03 = 6.9$ tolerances and a cost of about $48$. The alcohol-only point misses original extract by \SI{0.493}{\degree}P, which is $9.9$ tolerances and a cost of about $97$. The weighted point misses extract by \SI{0.167}{\degree}P and alcohol by \SI{0.137}{\percent} by volume, a cost of $(3.34)^{2}+(4.57)^{2}\approx 32$. The predictive controller is not applying a hidden preference for extract. It is applying the squares of the published tolerances. A brewery that considers a tenth of a percent of alcohol more serious than a tenth of a degree Plato changes two numbers in \eqref{eq:stage} and recommissions nothing else.

The flow weight is deliberately smaller than a quality tolerance. On this skid a shortfall of one cubic metre per hour is a scheduling inconvenience, and a shortfall of one degree Plato is a beer that cannot be released. If the skid were the filler and the shortfall were a missed production slot of a high-volume brand, the flow weight would rise. The formulation already has a place for that number.

\subsection{One shooting step, in order}
\label{sec:onestep}

A single evaluation of the horizon at the warm start, which is ``hold the lager command'', during the first sample of the light-brand grade change, does the following.

The mother-beer analysis is unchanged. The setpoint has become \SI{9.687}{\degree}P, \SI{4.000}{\percent} by volume, \SI{5.00}{\gram\per\litre} and \SI{20}{\cubic\metre\per\hour}. The mixer state is still the lager composition, because the valves have not moved. Both blocks of $\mathbf{z}$ are \eqref{eq:usteady}. The network solve returns the lager flows. The mixer step moves the state by $\alpha$ toward an inlet that is still the lager blend, so the predicted original extract stays near \SI{11.61}{\degree}P. The stage cost on extract alone is $((11.61-9.69)/0.05)^{2} \approx 1.5\times 10^{3}$ per step, and six steps cost about $10^{4}$. That number is the signal the solver follows. It is large because the beer in the vessel is still the previous brand, which is a fact, and because the command has not yet been asked to change, which is the degree of freedom.

The solver then shortens the beer lift and lengthens the water lift, subject to the rate limit of $0.16$ lift per four-second sample and to the slower rate inside the prediction. The new flows change $x_{\mathrm{in}}$ through \eqref{eq:dilute}. The predicted extract falls on the time constant \eqref{eq:tau}. Carbon dioxide is recomputed with the new water fraction and the new dose; the dose rises slightly because more water has to be carbonated even though the target concentration has fallen. The pressure residual is re-solved at every node so that the final valve opens as the total flow and the mix density change, instead of letting the pressure loop discover the interaction one sample later. The first block of the improved $\mathbf{z}$ is applied. The next sample repeats the procedure from the new mixer state and the new previous command.

Nothing in that sequence requires the outlet analyser. The analyser arrives \SI{12}{\second} later and trims whatever the flow integration got wrong. On the commissioned model that trim is at the level of the sensor noise. On the nominal model the trim cannot remove a bias that the model interprets as a correct steady state, which is the tail of \SI{0.19}{\degree}P in Section~\ref{sec:s1}.

\subsection{The production solver}
\label{sec:rti}

The case-study program is a single-shooting nonlinear program in eight variables with bound constraints and a smooth cost, apart from the hinges in the soft penalties. Those hinges are squared, so the cost is continuously differentiable and the second derivative jumps. A sequential quadratic solver handles that. The iteration limit of six is a real-time budget, not an optimality certificate. Warm-starting from the previous sample is what makes six iterations enough when the brand is not changing, and what makes the first sample of a grade change the expensive one.

The production form replaces those six iterations by one quadratic program per sample, which is the real-time iteration \citep{diehl2005,gros2020}. Linearise the shooting map at the previous solution $(\mathbf{z},\mathbf{x})$, build the quadratic model of \eqref{eq:stage}, solve it, and apply the first block before the next preparation phase has finished. The preparation of the derivatives can overlap the sample. IPOPT is a reliable way to solve the quadratic model and also the full nonlinear program during commissioning, when there is no sample clock \citep{wachter2006}. CasADi is a reliable way to obtain the derivatives of the network residual and of Balling's map without hand-written Jacobians \citep{andersson2019}. The thesis does not depend on either package being present in the case-study script. It depends on the shooting problem being the one those packages would be given.

Two implementation constraints follow. The scalar pressure solve has to be differentiated through, or replaced by a residual that the nonlinear program carries as an equality. Brent's method on a scalar is easy to run and awkward to differentiate. A production code lifts $p_m$ into the decision vector and writes $Q_f^{\mathrm{valve}}-Q_f^{\mathrm{mass}}=0$ as a constraint at each node. The local derivative of that residual with respect to $p_m$ is negative and bounded away from zero whenever both inlet valves have a positive drop, which they do inside the lift bounds of \eqref{eq:ineq}. The constraint qualification is then the ordinary one for a hydraulic network with a free internal pressure. The second constraint is that the rate limit and the interlocks are enforced in the programmable controller even if the quadratic program returns a point on the wrong side of them. Section~\ref{sec:twin} already assigned that duty to the PLC. The solver is not a safety system.

\subsection{Why roughness and minor loss share one column}
\label{sec:confound}

The line drop \eqref{eq:darcy} depends on $\varepsilon$ only through $f(\varepsilon)$ and on $K$ directly. At a fixed Reynolds number the map $(\varepsilon,K)\mapsto C_L$ is the only map the flows can see, because $C_L$ is the only combination that enters $\Delta p_{\mathrm{line}}$. Differentiating Haaland's formula, $f$ is a slowly increasing function of $\varepsilon/D$. A positive step in $\varepsilon$ and a negative step in $K$ can be chosen so that $C_L$ does not change. The two corresponding columns of the residual Jacobian are then scalar multiples of each other, the Gram matrix $J^{\mathsf T}J$ is singular, and the formal standard error \eqref{eq:cov} diverges. That is what Table~\ref{tab:est} records, in the decision to publish $C_L$ and to withhold a roughness.

The multiplicity is not removed by collecting more points at the same Reynolds number. It is weakly removed by collecting points at very different velocities, because $f$ depends on $\mathrm{Re}$ and $K$ does not, so the two parameters are not in a fixed ratio across the campaign. The commissioning lifts do change the flow, but not by the factor of ten that would separate a smooth-pipe friction factor from a minor loss. A campaign designed to identify roughness would include a near-shut final valve and a wide-open one, and would still be limited by the fact that a hygienic bore at these Reynolds numbers sits in a flat part of the Moody chart. The engineering quantity is $C_L$. The thesis estimates that quantity to two and a half percent and treats the split as unidentified.

The capacities do not have this problem. $K_{v,h}$ scales only $Q_h$, and it scales it with the square root of a drop that the header bias also shifts. Varying the beer lift changes the drop and the flow in a combination a pure bias cannot imitate, because a bias shifts every beer-valve point by the same pressure and a capacity scales them. The standard error of \SI{0.08}{\cubic\metre\per\hour} on a capacity of $32$ is the numerical expression of that excitation. Twenty-eight points were more than the rank of the problem required. They were not more than the noise required.

\subsection{The residual network as a function}
\label{sec:netfun}

Let $\xi\in\mathbb{R}^{6}$ be the standardised feature $(l_h,l_w,l_f,A_w,E_r,p_m)$, with the mean and the scale taken from the training points only. The network is
\begin{equation}
\label{eq:net}
h = \tanh(W_1\xi + b_1),
\qquad
\Delta Q = W_2 h + b_2,
\end{equation}
with $W_1\in\mathbb{R}^{8\times 6}$ and $W_2\in\mathbb{R}^{3\times 8}$. The corrected flow is the physical prediction plus $\Delta Q$. The training loss on a batch is the mean square of $\Delta Q$ against the measured residual, plus the weight decay, plus, every fifth step, a multiple of \eqref{eq:masspen}. The mass penalty is a function of $\Delta Q$ and therefore of $W_2$ and, through $h$, of $W_1$. Its gradient is an ordinary back-propagation of a linear form in the three flows. No collocation point in the sense of a space-time residual of a partial differential equation is required \citep{raissi2019}. The conservation law is the one the flow already lives under.

The hold-out is six points drawn as the last six of the campaign, not re-used in the standardisation or the training. Root-mean-square flow error rises from \SI{0.030}{\cubic\metre\per\hour} to \SI{0.037}{\cubic\metre\per\hour} when $\Delta Q$ is added. The rise is the signature of a model whose physical part has already reached the sensor noise. Adding eight hidden units, forty-eight input weights and twenty-four output weights to a residual whose standard deviation is \SI{0.03}{\cubic\metre\per\hour} on twenty-two points is enough capacity to fit the particular noise of those twenty-two points. The gate compares the two errors and publishes the physical prediction. The weights are kept on disk with the campaign so that a later campaign can be scored against the same architecture. They are not written into the controller.

This is the opposite deployment decision from \citet{wu2023cerd} and \citet{zheng2023}, and it is the same rule. They keep the network because the hold-out improves. This skid drops the network because the hold-out does not. A twin that kept the network anyway would be adding \SI{0.007}{\cubic\metre\per\hour} of invented flow, which is a third of a percent at the beer line and a pointless risk next to a capacity estimate that is already inside \SI{0.1}{\cubic\metre\per\hour}.

\subsection{Observer, delay and the information the analyser still carries}
\label{sec:observer}

The fast observer is the discrete mixer \eqref{eq:exactdisc} driven by measured flows. Its only unknown, between analyser updates, is the inlet composition, and that unknown is replaced by the last mother-beer analysis after eight seconds. The outlet analysis is a measurement of the state that existed \SI{12}{\second}, or six integration steps, earlier. A correction that simply overwrites the current state with the delayed analysis assigns today's vessel the beer of one residence time ago, which is the lager during the first moments of the grade change, and pulls the prediction backwards. The correction used here is therefore applied to the innovation of the delayed prediction: the observer stores a short buffer of predicted outlet values, subtracts the analysis when it arrives, and distributes a fraction of that innovation onto the current state. With healthy flow meters the fraction can be small. The case study adds composition noise of about $0.01$ percentage points by weight so that the controller does not receive the plant state as an oracle.

If the beer flow meter failed, this observer would integrate a lie and the outlet analysis would become the primary measurement. That is the situation in which the moving-horizon estimator of \citet{rao2003} earns its window, and in which an unscented filter \citep{julier2004} is a cheaper approximation. The thesis does not simulate a failed meter. It names the estimator that would replace the balance, and it keeps the balance for the campaign that was actually run. \citet{kalman1960} is the linear special case of both, and it is the right filter for the pressure and flow channels, which are already linear in their noise and are used without a window.

Dissolved oxygen is observed on the same analyser and is not controlled. The observer reports it. The phase logic compares it with \SI{30}{ppb}. The predictive controller includes it in $\ell_{\mathrm{soft}}$ so that a move which would add oxygen by flashing is expensive, and so that a move which cannot remove oxygen already present is not invented. Those are different uses of the same measurement. Conflating them is how a ratio station gets talked into chasing a disturbance it has no actuator for.

\subsection{Faults the cost already contains}
\label{sec:faults}

Four faults are inside the formulation rather than on a list of future work.

A biased beer-header transmitter is $b_h$. It is estimated once and subtracted. Leaving it inside the nominal model produces the grade-change offset of Section~\ref{sec:s1}.

A wrong capacity is an element of $\theta$. It is the difference between the two predictive controllers in every scenario. The ratio station is less sensitive to it because the flow loops close on measured flow. The predictive controller is more sensitive to it because the flow it closes on is a prediction. That asymmetry is the reason the twin is not an optional display next to an otherwise finished controller.

An infeasible brand is a positive minimum of \eqref{eq:stage}. The controller returns the minimiser. It does not return a feasibility flag unless the supervisory station computes the same one-dimensional search in $w$ that Section~\ref{sec:s2} used and compares the two pure compromises with the weighted point. That search is cheap and should be part of the phase start, so that an operator sees ``this mother beer cannot hit both numbers'' before the valves move. The case study performs the search offline and confirms that the closed-loop tail matches it.

A sticky water valve is $S$ in \eqref{eq:invstick}. Undetected, it turns a small flow correction into a delayed jump. Detected, it is a known lead on one command. The eight-second detection delay in the third scenario is a stand-in for a test that looks for a command change with no stem change \citep{jelali2010,choudhury2005}. It is not a claim that eight seconds is the best detector.

Faults that are not inside the formulation are a plugged injector, a failed analyser, a bright-tank pressure below $p_{\mathrm{eq}}$, and a dilution water whose calcium or oxygen certificate is already outside the release limit. The first three are interlocks. The fourth is a hold on the phase. The controller is not asked to blend its way around them.


\subsection{What each technology is for}

The first technology is a way to choose $\mathbf{u}$ when the model is believed. It coordinates the three valves, it uses the nonlinear quality map, and it prices a flashing vena contracta in the same sum as a missed extract. The second technology is a way to decide whether the model should be believed, and to correct the particular parameters that a blend skid gets wrong: capacities, a pressure bias, a solubility scale, a density contraction, a line resistance that must be estimated as a sum, and a deadband. A residual network is part of the second technology and is subject to a hold-out gate.

The next chapter builds the numerical skid on which both are estimated, and it reports what the beer does.
